\documentclass[11pt,a4paper]{article}
\usepackage[spanish]{babel}
\usepackage[T1]{fontenc}
\usepackage[utf8]{inputenc}
\usepackage{amsmath,booktabs,geometry,hyperref}
\geometry{margin=25mm}
\title{S01: reporte de un experimento de modelos probabilísticos de texto}
\author{Nombre del estudiante: \rule{65mm}{.3pt}}
\date{Versión del código y fecha: \rule{45mm}{.3pt}}
\begin{document}\maketitle
\section{Pregunta e hipótesis}
Pregunta: ¿un bigrama entrenado en textos técnicos fabricados predice las
seis frases técnicas reservadas con menor BPC que otro entrenado en
prosa general fabricada? Declare la dirección del contraste.
\section{Datos y licencias}
El conjunto \texttt{corpus\_didactico.csv} contiene textos sintéticos
creados para la clase, no informes reales. Ocho y ocho para ajustar; seis
para evaluar. Indique cómo se evitó usar prueba en entrenamiento.
\section{Métodos y reproducibilidad}
Bigrama a nivel de carácter, $\alpha=1$, vocabulario común fijado en
entrenamiento, BOS/EOS, Unicode como puntos de código, minúsculas,
recorte de espacios de extremos. Nombre comandos y versiones ejecutadas.
\section{Resultados (debe regenerarlos)}
\begin{center}\begin{tabular}{lcc}\toprule
Estimador & Valor de ejemplo & Interpretación \\\midrule
Diferencia media, BPC (técnico - general) & $-0.54024842$ & Solo corpus sintético \\
IC bootstrap percentil pareado (95\%) & $[-0.66771,-0.42038]$ & Remuestreo de seis frases \\\bottomrule
\end{tabular}\end{center}
Reemplace o confirme los valores con su ejecución. Incluya la tabla de
resultados por documento y la versión del intérprete. Si R no está
instalado, declare la ejecución pendiente.
\section{Microexperimento IC10: dos tokenizadores BPE}
Comparamos BPE general y tecnico con \textbf{20 fusiones} cada uno,
entrenados exclusivamente con sus corpus respectivos y evaluados en
las mismas seis frases tecnicas reservadas. Son 35 palabras en total.
\begin{center}\begin{tabular}{lrr}\toprule
& Tokens & Tokens por palabra\\\midrule
BPE general &148&4.228571\\
BPE tecnico &127&3.628571\\\bottomrule
\end{tabular}\end{center}
Incluya \texttt{resultados\_ic10\_tokenizacion.csv}, su comando de
ejecucion y los datos por documento. Estas cifras describen longitud
tokenizada, \emph{no} perdida predictiva ni resultados de LLM.
\section{Amenazas a la validez y conclusión responsable}
No hay muestra representativa ni comparacion con tokenizadores modernos;
la reduccion de longitud en el microexperimento no implica mejora de BPC.
No existen evidencias generalizables sobre una poblacion de reportes. Explique por qué el objetivo de un
artículo Q1 exige novedad, baselines fuertes y evaluación externa.
\section{Puente con investigación: dos áreas de interés}
\paragraph{Área 1.} Relevancia, datos candidatos y riesgo.
\paragraph{Área 2.} Relevancia, datos candidatos y riesgo.
\end{document}
